\section{Correctness Analysis of the Reverse-Cuckoo Evaluator}
\label{app:revCuckooHash}

This appendix proves \namedref{Proposition}{prop:rc-ringlpn-correctness}.  The
analysis is about consistency of the programmed binary systems.  It is
independent of the Reverse-Cuckoo leakage assumption in
\sectionref{sec:revCuckooDistribution}.

\subsection{Caller distribution and shuffle factors}

Use the notation of \equationref{eq:ringlpn-product-list}.  At the claimed
Ring-LPN point,
\[
 N_{\mathsf R}=2^{20},
 \qquad P=4,
 \qquad t=16,
 \qquad L=N_{\mathsf R}/t=2^{16}.
\]
The $P^2t=256$ triples $(a,b,k)$ produce 256 Reverse-Cuckoo sets.  Before
deduplication, set $A_{a,b,k}$ contains the $t$ sums
\[
 X_{0,a,j}+X_{1,b,(k-j)\bmod t},
 \qquad j\in[0,t).
\]
Before any support rejection, these points within one set are independent samples from the triangular
distribution obtained by adding two uniform elements of $[0,L)$.  The DPF
domain is $[0,2L)$, which we identify through binary encoding with
$\mathcal X:=\F_2^{17}$.  Every point has probability at most
$2/|\mathcal X|$.  Different sets share the underlying offsets and are
therefore correlated.  The union bound below uses only expected relation
counts; it does not assume independent failures across sets.

Fix a relation of weight $k$.  Under a uniform shuffle into $w$ partitions of
$d$ rows, the probability that all its points enter one partition is
\begin{equation}\label{eq:rc-one-partition-prob}
 p_{\mathsf{part}}(k)=\frac{w(d)_k}{(wd)_k},
\end{equation}
where $(a)_k=a(a-1)\cdots(a-k+1)$.  Conditioned on this event, its labels are a
uniform $k$-subset of $\F_2^{\log_2d}$.  Let $p_0(d,k)$ be the probability that
their XOR is zero.  A character sum gives
\begin{equation}\label{eq:rc-zero-label-prob}
 p_0(d,k)=
 \begin{cases}
  \displaystyle\frac1d, & k\text{ odd},\\[2mm]
  \displaystyle\frac{\binom dk+(d-1)(-1)^{k/2}\binom{d/2}{k/2}}
                         {d\binom dk}, & k\text{ even}.
 \end{cases}
\end{equation}
The relation is harmful with the remaining probability $1-p_0(d,k)$.

\subsection{Counting all degree-two relations}

Only relations with $L_\alpha=0$ can disappear in the nonlinear lift.  For
real points, the sentinel coordinate also implies that such a relation has
even weight.  Let $B[k,s]$ denote the number of $k$-subsets of $\mathcal X$
whose XOR is zero and whose quadratic moment has rank $2s$.  The rank-zero
class consists of the weight-$k$ words in $\operatorname{RM}(2,17)^\perp$;
the other classes retain nonzero quadratic moments.

The refinement can be computed without enumerating point sets.  Let
$K=\dim\operatorname{RM}(2,17)=1+17+\binom{17}{2}$, let $v_i$ be the number of
alternating forms of rank $2i$, and let $P_{i,s}$ be the eigenmatrix of the
binary alternating-forms association scheme.  If $T_{k,i}$ is the sum of the
weight-$k$ Hamming Krawtchouk polynomial over the affine parts of a quadratic
form of rank $2i$, then
\begin{equation}\label{eq:rc-rank-enumerator}
 B[k,s]=2^{-K}\sum_i v_iP_{i,s}T_{k,i}.
\end{equation}
We use the alternating-forms eigenmatrix of Delsarte and
Goethals~\cite{delsarteGoethals1975}.  \equationref{eq:rc-rank-enumerator}
specializes the MacWilliams transform while retaining quadratic rank; direct
enumeration in dimensions three and four verifies the implementation of the
transform.

For any fixed $k$-subset, the probability that all its points occur among
the $t$ samples is at most $(t)_k(2/|\mathcal X|)^k$.  Deduplication cannot increase the
number of distinct relations.  The expected number of weight-$k$, rank-$2s$
relations in the complete batch is therefore at most
\begin{equation}\label{eq:rc-expected-relation-count}
 P^2t\,(t)_k B[k,s]\left(\frac{2}{|\mathcal X|}\right)^k.
\end{equation}
For such a relation, the probability that the degree-two lift vanishes is
\begin{equation}\label{eq:rc-lift-collapse}
 a_{q'}(2s):=\left(\frac12+2^{-2s-1}\right)^{q'}.
\end{equation}
Combining \namedeqref{Equations}{eq:rc-one-partition-prob},
\eqref{eq:rc-zero-label-prob}, \eqref{eq:rc-expected-relation-count}, and
\eqref{eq:rc-lift-collapse}, then summing over $k\leq t$ and $s$, gives the
structural failure bound
\begin{equation}\label{eq:rc-structural-bound}
 \epsilon_{\mathsf{str}}\leq
 \sum_{k,s} P^2t\,(t)_k B[k,s]
 \left(\frac{2}{|\mathcal X|}\right)^k
 \frac{w(d)_k}{(wd)_k}
 \bigl(1-p_0(d,k)\bigr)a_{q'}(2s).
\end{equation}
This sum includes deterministic degree-two identities: when $s=0$,
$a_{q'}(0)=1$.

For the numerical bounds, we sharpen the expected count of affine planes
using the triangular distribution directly.  Write $p(x)$ for its mass at
$x\in\mathcal X$, and define its Walsh transform by
$\widehat p(u):=\sum_x p(x)(-1)^{u\cdot x}$.  The probability that four
independent samples are distinct and have XOR zero is
\[
 \frac{1}{|\mathcal X|}\sum_u\widehat p(u)^4
 -3\left(\sum_x p(x)^2\right)^2+2\sum_x p(x)^4.
\]
Multiplying by $P^2t\binom t4$ bounds their expected count.  The numerical
structural sum also conservatively uses the full survival probability in
\equationref{eq:rc-relation-survival} instead of only $a_{q'}(2s)$.
This counts some final-compression events twice but preserves an upper bound.

The largest terms for $w=2$, $d=16$, and the former intermediate width
$q'=56$ are shown below.  Every omitted rank class contributes less than the
last displayed row.
\begin{table}[H]
\centering
\caption{Dominant structural terms at $q'=56$.}
\label{tab:rc-hash-structural-terms}
\begin{tabular}{c@{\qquad}c@{\qquad}c}
\toprule
Relation weight $k$ & Quadratic rank $2s$ & Batch contribution \\
\midrule
$4$  & $2$ & $2^{-39.469}$ \\
$8$  & $0$ & $2^{-42.217}$ \\
$6$  & $4$ & $2^{-46.993}$ \\
$8$  & $6$ & $2^{-50.366}$ \\
$10$ & $8$ & $2^{-52.928}$ \\
\bottomrule
\end{tabular}
\end{table}
The first row consists of affine planes.  A separate Walsh-transform
calculation bounds their expected number before the shuffle by $3.791522$ and
gives the same $2^{-39.469}$ contribution.  Rank-zero relations begin at
weight eight and are affine three-cubes; their complete contribution is at
most $2^{-42.217}$.  Thus deterministic degree-two identities are included,
not excluded by a sampling heuristic.

\subsection{Final compression and complete bound}

Fix the supports, shuffle, and lifted rows in one partition containing $r$
real rows.  Each nonzero lifted relation is killed by the independent uniform
matrix $A_2$ with probability $2^{-q}$.  There are at most $2^r-1$ such
relations, so their union contribution is less than $2^{r-q}$.  This argument
does not require the lifted rows to be independent.

Under a uniform shuffle, the number of real rows in one partition is
hypergeometric before deduplication.  Deduplication can only reduce this
number.  Averaging and taking a union bound over all sets and partitions gives
\begin{equation}\label{eq:rc-compression-bound}
 \epsilon_{\mathsf{cmp}}
 <P^2tw\sum_{r=0}^{t}
 \frac{\binom dr\binom{(w-1)d}{t-r}}{\binom{wd}{t}}\,2^{r-q}.
\end{equation}
The structural and compression events cover every nonzero row relation:
either it vanishes in the lift, or a nonzero lifted relation is killed by
$A_2$.  Hence
$\epsilon_{\mathsf{batch}}\leq
\epsilon_{\mathsf{str}}+\epsilon_{\mathsf{cmp}}$.

\begin{table}[H]
\centering
\caption{Analytic correctness bounds before support rejection for the
$N_{\mathsf R}=2^{20}$, $w=2$, $d=16$ Ring-LPN batch.
Displayed exponents are rounded.}
\label{tab:rc-hash-complete-bound}
\begin{tabular}{c@{\qquad}c@{\qquad}c@{\qquad}c}
\toprule
$q'=q$ & Structural & Compression & Complete \\
\midrule
$56$ & $2^{-39.261}$ & $2^{-38.288}$ & $2^{-37.694}$ \\
$64$ & $2^{-42.007}$ & $2^{-46.288}$ & $2^{-41.935}$ \\
\bottomrule
\end{tabular}
\end{table}
Substituting $d=16$ and $\lambda=40$ in
\equationref{eq:rc-hash-widths} gives $q'=q=64$, proving the unfiltered
clause of \namedref{Proposition}{prop:rc-ringlpn-correctness}.  For comparison, the same
calculation for $w=3$, $d=8$ gives only $32.577$ bits at width $48$, but
$40.194$ bits at the fixed-margin width $56$.

The relation count also depends on the ring degree because
$L=N_{\mathsf R}/t$ determines the point domain.  Applying the same analytic
calculation to the three Ring-LPN benchmark sizes gives the following bounds.
\begin{table}[H]
\centering
\caption{Dependence of the $q'=q=64$ proof on the Ring-LPN ring degree.}
\label{tab:rc-hash-ring-size}
\begin{tabular}{c@{\qquad}c@{\qquad}c@{\qquad}c}
\toprule
$N_{\mathsf R}$ & $L=N_{\mathsf R}/t$ & DPF domain & Complete bound \\
\midrule
$2^{16}$ & $2^{12}$ & $2^{13}$ & $2^{-26.218}$ \\
$2^{18}$ & $2^{14}$ & $2^{15}$ & $2^{-34.213}$ \\
$2^{20}$ & $2^{16}$ & $2^{17}$ & $2^{-41.935}$ \\
\bottomrule
\end{tabular}
\end{table}
At the two smaller ring degrees, rank-zero affine cubes dominate the bound.
Their relations vanish in every degree-two feature lift.  Increasing the two
evaluator widths therefore does not certify 40-bit correctness at those
points; a stronger evaluator or a sharper caller-specific argument would be
required.

The proof assumes a uniform shuffle, ideal random matrices derived from the
seed, uniform Ring-LPN block offsets, and semi-honest execution.  It does not
analyze malicious bias in joint seed generation.  It also does not apply to
the AnyField caller, whose point distribution must be counted separately.

\subsection{Conditioning on the folded-support filter}
\label{app:rc-hash-conditioning}

The weak-factor filter changes the support distribution, so the preceding
bound cannot be applied unchanged.  Each party independently resamples its
four supports until their total folded support weight is at least 61.
The weight counts occupied residues separately in each polynomial, before
coefficient cancellation.  For party $p\in\{0,1\}$ and polynomial $a\in[0,P)$,
define
\[
 D_{p,a}:=\left|\{(jL+X_{p,a,j})\bmod128:j\in[0,t)\}\right|,
 \qquad
 R_p:=\left\{\sum_{a\in[0,P)}D_{p,a}\ge61\right\}.
\]
The shuffle and evaluator seeds remain independent of this filter.
Here $128$ divides $L$, so before rejection the $D_{p,a}$ are independent
copies of the number $D$ of occupied bins after 16 uniform draws into 128 bins.

Let $D_1,\ldots,D_4$ be independent copies of $D$, and put
\[
 \alpha:=\Pr[D_1+\cdots+D_4\ge61],
 \qquad
 \beta:=\Pr[D_1+D_2+D_3\ge45].
\]
These probabilities follow from an occupancy recurrence.  If $u_j(v)$ is
the probability of $v$ occupied bins after $j$ draws, then $u_0(0)=1$ and
\[
 u_{j+1}(v)=\frac{v}{128}u_j(v)
          +\frac{129-v}{128}u_j(v-1),
\]
with out-of-range entries zero.  Convolving $u_{16}$ gives
$\alpha\approx0.5019144583$ and $\beta\approx0.7174085984$.

Fix one polynomial's complete offset vector.  It has at most 16 occupied
residues, so the probability that the other three polynomials complete
acceptance is at most $\beta$.  Bayes' rule therefore bounds the density of
this polynomial after local rejection by $\beta/\alpha$ times its original
density.  The two parties reject independently.  Any nonnegative function
of one polynomial from each party and independent setup coins consequently
satisfies
\begin{equation}\label{eq:rc-pair-conditioning}
 \mathbb E[F\mid R_0\cap R_1]
 \le \left(\frac{\beta}{\alpha}\right)^2\mathbb E[F].
\end{equation}

Every product list $A_{a,b,k}$ depends on just such a pair of polynomials.
Apply \equationref{eq:rc-pair-conditioning} to each list's structural
union contribution and sum over all lists.  This does not assume that
different lists fail independently.  The final-compression bound is
pointwise in the supports and therefore remains unchanged.  We obtain
\begin{equation}\label{eq:rc-filtered-hash-bound}
 \epsilon_{\mathsf{batch}\mid R_0\cap R_1}
 \le \left(\frac{\beta}{\alpha}\right)^2\epsilon_{\mathsf{str}}
     +\epsilon_{\mathsf{cmp}}
 <4.739\cdot10^{-13}<2^{-40.940}.
\end{equation}
Here $(\beta/\alpha)^2\approx2.043025290$.  Thus the proposed support-only
filter preserves 40-bit programmed-system correctness at the stated
$2^{20}$ ring degree, without changing the evaluator widths.  This is a
correctness bound, not a bound on descriptor leakage.  A filter that also
tests coefficients, or a different rejection rule, requires its own analysis.
